% Complete Section 4.5 and Box 4.1, original PDF 54-60 / printed 44-50.
% The printed signs and the last summand in Box 4.1 are retained, not corrected.
% block: hodge-r3-heading page:54 type:section
\subsection{{Example: The Hodge Laplacian in $\mathbb R^3$}}\label{sec:hodge-r3}

% block: hodge-r3-introduction pages:54,55 type:ordinary-paragraph
{We now return to our running example, the de Rham complex on a domain with Lipschitz boundary in $\mathbb R^3$. We saw in Section~\ref{sec:l2-de-rham-complex} that this forms a closed Hilbert complex. Therefore, for each $k=0,1,2,3$, we obtain a Hodge Laplace problem associated to the complex, which is a well-posed boundary value problem in partial differential equations. To interpret these problems we just have to interpret the spaces $V^k$ and $V_k^*$ and the operators $d^k$ and $d_k^*$, for which we will have frequent recourse to the domain complex \eqref{eq:l2-de-rham-domain-complex} and the dual domain complex \eqref{eq:l2-dual-domain-complex}.}

% block: hodge-r3-k0-heading page:55 type:subsection
\subsubsection{$k=0$}\label{subsec:hodge-r3-k0}

% block: hodge-r3-k0-domain page:55 type:ordinary-paragraph
{Assuming that the domain is connected, the space $\mathsf H^0$ of harmonic forms consists only of the constants and $P_{\mathsf H}f=\bar f$, the mean value of $f$ on $\Omega$. Since $d_0^*=0$, $L^0u=(dd^*+d^*d)u=d^*du=-\operatorname{div}\operatorname{grad}u=-\Delta u$, with domain}
\[
\begin{aligned}
D(L)&=\{\,u\in H^1(\Omega)\mid\operatorname{grad}u\in\mathring H(\operatorname{div},\Omega)\,\}\\
&=\{\,u\in H^1(\Omega)\mid\Delta u\in L^2(\Omega),\ \partial u/\partial n=0\text{ on }\partial\Omega\,\},
\end{aligned}
\]
{where $\Delta$ denotes the ordinary Laplacian $\Delta=\sum_{i=1}^3\partial^2/\partial x_i^2$. Given $f\in L^2(\Omega)$, the strong formulation \eqref{eq:hodge-strong-formulation} then seeks $u\in H^1$ such that}
\begin{equation}\label{eq:hodge-r3-neumann-strong}
-\Delta u=f-\bar f\text{ in }\Omega,\qquad\frac{\partial u}{\partial n}=0\text{ on }\partial\Omega,\qquad\int u=0.
\end{equation}

% block: hodge-r3-neumann-interpretation page:55 type:ordinary-paragraph
{Thus the $k=0$ Hodge Laplace problem associated to the de Rham complex is just the usual Neumann problem for Poisson's equation, the adjustment by $\bar f$ being required to make the data $f$ compatible with the homogeneous Neumann boundary condition.}

% block: hodge-r3-neumann-weak page:55 type:ordinary-paragraph
{Next we turn to the primal and mixed weak formulations. Since $V^{-1}=0$, for $k=0$, $\sigma$ and $\tau$ vanish in the mixed weak formulation \eqref{eq:hodge-mixed-formulation}. Thus there is no distinction between the two weak formulations for $k=0$. They both seek $u\in H^1$ such that}
\begin{equation}\label{eq:hodge-r3-neumann-weak}
\int\operatorname{grad}u\cdot\operatorname{grad}v=\int(f-\bar f)v,\qquad v\in H^1,\qquad\int_\Omega u=0.
\end{equation}

% block: hodge-r3-neumann-equivalence page:55 type:ordinary-paragraph
{We have seen in Theorem~\ref{thm:hodge-formulation-equivalence} that this problem is equivalent to the strongly formulated problem \eqref{eq:hodge-r3-neumann-strong}. For concreteness, we now check directly that \eqref{eq:hodge-r3-neumann-weak} implies \eqref{eq:hodge-r3-neumann-strong}, the other direction being easier. Taking $v\in C_0^\infty(\Omega)$ in \eqref{eq:hodge-r3-neumann-weak} we see that}
\begin{equation}\label{eq:hodge-r3-neumann-distribution}
-\operatorname{div}\operatorname{grad}u=f-\bar f
\end{equation}
{in the sense of distributions and so, since $f\in L^2$, $\operatorname{grad}u$ belongs to $H(\operatorname{div})$. We may therefore apply the integration-by-parts formula \eqref{eq:normal-trace-green} with $w=\operatorname{grad}u$ and then apply \eqref{eq:hodge-r3-neumann-weak} and \eqref{eq:hodge-r3-neumann-distribution} to get}
\[
\langle\gamma_n(\operatorname{grad}u),\operatorname{tr}v\rangle=\int(\operatorname{div}\operatorname{grad}u)v+\int\operatorname{grad}u\cdot\operatorname{grad}v=0,\quad v\in H^1.
\]
{Since $\operatorname{tr}v$ can be an arbitrary element of $H^{1/2}(\partial\Omega)$ we conclude that $\gamma_n(\operatorname{grad}u)=0$, i.e., $\partial u/\partial n=0$ on the boundary. Notice that, in the weak formulation, the boundary condition is imposed \emph{naturally}, meaning that it is implied by the equations via integration by parts rather than being imposed in the space in which $u$ is sought.}

% block: hodge-r3-k0-boundary-problems page:55 type:ordinary-paragraph
{Turning to the $\mathcal B$ and $\mathcal B^*$ problems, we do not find anything new, since $\mathcal B^0=0$. Therefore, the $\mathcal B$ problem does not arise for $k=0$, while the $\mathcal B^*$ problem is essentially the full Hodge Laplacian.}

% block: hodge-r3-k0-summary page:56 type:ordinary-paragraph
{Summarizing the situation for $k=0$, the Hodge Laplacian for $0$-forms is just the Neumann problem for the Poisson equation. The primal and mixed weak formulations coincide in this case, being just the standard weak formulation for the Neumann problem, in which the boundary condition is imposed \emph{naturally}.}

% block: hodge-r3-k1-heading page:56 type:subsection
\subsubsection{$k=1$}\label{subsec:hodge-r3-k1}

% block: hodge-r3-k1-strong page:56 type:ordinary-paragraph
{The case $k=1$ is more interesting. The space of harmonic forms $\mathsf H^1$ is a space of vector fields whose dimension is $b_1$, the first Betti number. The strong formulation seeks $u\in H(\operatorname{curl})\cap\mathring H(\operatorname{div})$ with $\operatorname{curl}u\in\mathring H(\operatorname{curl})$ and $\operatorname{div}u\in H^1$ satisfying the partial differential equation}
\begin{equation}\label{eq:hodge-r3-vector-laplace}
\operatorname{curl}\operatorname{curl}u-\operatorname{grad}\operatorname{div}u=f-P_{\mathsf H}f\text{ in }\Omega.
\end{equation}
{Note that the differential operator $\operatorname{curl}\operatorname{curl}-\operatorname{grad}\operatorname{div}$ is the negative vector Laplacian $-\vec\Delta$ in that \[\operatorname{curl}\operatorname{curl}u-\operatorname{grad}\operatorname{div}u=(-\Delta u_1,-\Delta u_2,-\Delta u_3)\] for a smooth vector field $u=(u_1,u_2,u_3)$.}

% block: hodge-r3-magnetic-boundary page:56 type:ordinary-paragraph
{The definition of the domain $D(L^1)$ implies the boundary conditions}
\begin{equation}\label{eq:hodge-r3-magnetic-boundary}
u\cdot n=0,\qquad(\operatorname{curl}u)\times n=0\text{ on }\partial\Omega,
\end{equation}
{and, of course, we have the side condition $u\perp\mathsf H^1$ to impose uniqueness. The conditions \eqref{eq:hodge-r3-magnetic-boundary} are sometimes referred to as \emph{magnetic boundary conditions}, for reasons that will become clearer when we discuss Maxwell's equations in Section~\ref{sec:maxwell-equations}.}

% block: hodge-r3-k1-primal page:56 type:ordinary-paragraph
{The primal weak formulation seeks $u\in H(\operatorname{curl})\cap\mathring H(\operatorname{div})$ with}
\begin{equation}\label{eq:hodge-r3-magnetic-primal}
\begin{gathered}
\langle\operatorname{curl}u,\operatorname{curl}v\rangle+\langle\operatorname{div}u,\operatorname{div}v\rangle=\langle f-P_{\mathsf H}f,v\rangle,\\
v\in H(\operatorname{curl})\cap\mathring H(\operatorname{div}).
\end{gathered}
\end{equation}
{Notice that, in this formulation, the boundary condition $u\cdot n=0$ is essential (imposed in the space where $u$ is sought), while the condition $\operatorname{curl}u\times n=0$ is natural.}

% block: hodge-r3-k1-mixed page:56 type:ordinary-paragraph
{By contrast, the mixed weak formulation seeks $\sigma\in H^1$, $u\in H(\operatorname{curl})$, $p\in\mathsf H^1$, without any imposed boundary conditions, satisfying}
\begin{equation}\label{eq:hodge-r3-magnetic-mixed}
\begin{aligned}
\int\sigma\tau-\int u\cdot\operatorname{grad}\tau&=0,&\tau&\in H^1,\\
\int\operatorname{grad}\sigma\cdot v+\int\operatorname{curl}u\cdot\operatorname{curl}v+\int p\cdot v&=\int f\cdot v,&v&\in H(\operatorname{curl}),\\
\int u\cdot q&=0,&q&\in\mathsf H^1.
\end{aligned}
\end{equation}
{The boundary conditions \eqref{eq:hodge-r3-magnetic-boundary} are both natural in this formulation, arising from the integration by parts in the first and second equations, respectively. Notice that, in addition to the unknown vector field $u$, this formulation takes as a second unknown the scalar function $\sigma=-\operatorname{div}u$ (and of course $p=P_{\mathsf H}f$).}

% block: hodge-r3-k1-b-problem page:56 type:ordinary-paragraph
{In the case $f\in\mathcal B^1$, i.e., $f$ is a gradient, then the solution $u$ is a gradient as well, which we can express by saying that $\operatorname{curl}u=0$ and $u\perp\mathsf H^1$. Thus the $\mathcal B^1$ problem is to find $u$ such that}
\[
-\operatorname{grad}\operatorname{div}u=f,\ \operatorname{curl}u=0\text{ in }\Omega,\qquad u\cdot n=0\text{ on }\partial\Omega,\qquad u\perp\mathsf H^1.
\]

% block: hodge-r3-div-curl-application page:57 type:ordinary-paragraph
{From the abstract theory, we know that this problem has a unique solution for $f\in\mathcal B^1$ which, moreover, can be computed by solving the full Hodge Laplacian problem for that same $f$ (e.g., via the mixed weak formulation). One application of such a problem is toward solving the div-curl problem}
\[
\operatorname{div}u=F,\qquad\operatorname{curl}u=0,
\]
{which can be done by taking $f=\operatorname{grad}F$ (with $\int F=0$ required for compatibility with the boundary condition on $u$).}

% block: hodge-r3-k1-bstar-problem page:57 type:ordinary-paragraph
{On the other hand, the $\mathcal B^*$ problem arises when $f\in\mathcal B_1^*$, meaning $f\in\operatorname{curl}\mathring H(\operatorname{curl})$. Then $u$ is determined by}
\[
\begin{gathered}
\operatorname{curl}\operatorname{curl}u=f,\ \operatorname{div}u=0\text{ in }\Omega,\\
u\cdot n=0,\ \operatorname{curl}u\times n=0\text{ on }\partial\Omega,\qquad u\perp\mathsf H^1.
\end{gathered}
\]
{Such curl-curl problems arise in electromagnetics. Again, the solution may be determined by solving the full Hodge Laplace problem.}

% block: hodge-r3-k2-heading page:57 type:subsection
\subsubsection{$k=2$}\label{subsec:hodge-r3-k2}

% block: hodge-r3-k2-domain page:57 type:ordinary-paragraph
{This case again leads to the vector Laplace equation \eqref{eq:hodge-r3-vector-laplace}, but with a different set of boundary conditions (and different harmonic forms with $\dim\mathsf H^2=b_2$). The domain of $L^2$ is the set of $u\in H(\operatorname{div})\cap\mathring H(\operatorname{curl})$ with $\operatorname{div}u\in\mathring H^1$ and $\operatorname{curl}u\in H(\operatorname{curl})$. Thus we obtain \emph{electric boundary conditions}:}
\[
u\times n=0,\ \operatorname{div}u=0\text{ on }\partial\Omega.
\]

% block: hodge-r3-k2-mixed page:57 type:ordinary-paragraph
{The first of these conditions is essential in the primal variational formulation, while the second is natural. Both are, of course, natural in the mixed weak formulation, which seeks $\sigma\in H(\operatorname{curl})$, $u\in H(\operatorname{div})$, $p\in\mathsf H^2$ such that}
\begin{equation}\label{eq:hodge-r3-electric-mixed}
\begin{aligned}
\int\sigma\cdot\tau-\int u\cdot\operatorname{curl}\tau&=0,&\tau&\in H(\operatorname{curl}),\\
\int\operatorname{curl}\sigma\cdot v+\int\operatorname{div}u\operatorname{div}v+\int p\cdot v&=\int f\cdot v,&v&\in H(\operatorname{div}),\\
\int u\cdot q&=0,&q&\in\mathsf H^2.
\end{aligned}
\end{equation}
{Notice that this formulation is very different from the mixed weak formulation \eqref{eq:hodge-r3-magnetic-mixed} for $k=1$. Here the additional field $\sigma$ is a vector field, while it was a scalar field for $k=1$.}

% block: hodge-r3-k2-boundary-problems page:57 type:ordinary-paragraph
{The $\mathcal B$ problem for $k=2$ is}
\[
\operatorname{curl}\operatorname{curl}u=f,\ \operatorname{div}u=0\text{ in }\Omega,\qquad u\times n=0,\text{ on }\partial\Omega,\qquad u\perp\mathsf H^2,
\]
{which has a unique solution if $f\in\operatorname{curl}H(\operatorname{curl})$. The $\mathcal B_2^*$ problem is}
\[
\begin{gathered}
-\operatorname{grad}\operatorname{div}u=f,\ \operatorname{curl}u=0\text{ in }\Omega,\\
u\times n=0,\ \operatorname{div}u=0\text{ on }\partial\Omega,\qquad u\perp\mathsf H^2,
\end{gathered}
\]
{solvable for $f\in\operatorname{grad}\mathring H^1$. Notice that the differential equations, but not the boundary conditions, for the $\mathcal B^2$ problem coincide with those of the $\mathcal B_1^*$ problem and those for the $\mathcal B_2^*$ problem coincide with those of the $\mathcal B^1$ problem.}

% block: hodge-r3-table page:58 type:table
\begin{table}[htbp]\centering\small
\caption{{The Hodge Laplace problem for the de Rham complex. $\sigma$ is the extra variable in mixed weak formulation. Underlined boundary conditions (the ones not involving derivatives) are essential in the primal weak formulation, while the others are natural. All the boundary conditions are natural in the mixed weak formulation.}}\label{tab:hodge-r3-boundary-conditions}
\setlength{\tabcolsep}{4pt}
\begin{tabular}{cccll}\hline
$k$&$L^k$&$\sigma$&{Boundary conditions}&{Harmonic forms}\\\hline
$0$&$-\Delta$&$0$&Neumann: $\partial u/\partial n=0$&$\mathsf H^0$ $(\dim b_0)$\\\noalign{\vskip4pt}
$1$&$-\vec\Delta$&$-\operatorname{div}u$&{magnetic}: $\underline{u\cdot n=0}$, $\operatorname{curl}u\times n=0$&$\mathsf H^1$ $(\dim b_1)$\\\noalign{\vskip4pt}
$2$&$-\vec\Delta$&$\operatorname{curl}u$&{electric}: $\underline{u\times n=0}$, $\operatorname{div}u=0$&$\mathsf H^2$ $(\dim b_2)$\\\noalign{\vskip4pt}
$3$&$-\Delta$&$-\operatorname{grad}u$&Dirichlet: $\underline{u=0}$&$\mathsf H^3=0$\\\hline
\end{tabular}
\end{table}

% block: hodge-r3-k3-heading page:58 type:subsection
\subsubsection{$k=3$}\label{subsec:hodge-r3-k3}

% block: hodge-r3-k3-dirichlet page:58 type:ordinary-paragraph
{Both $\mathcal B_3^*$ and $\mathsf H^3$ vanish. The domain of the Laplacian $D(L^3)=\{\,u\in\mathring H^1\mid\Delta u\in L^2\,\}$, so the boundary value problem is just the Dirichlet problem for Poisson's equation:}
\[
-\Delta u=f\text{ in }\Omega,\qquad u=0\text{ on }\partial\Omega.
\]
{The primal weak formulation is the obvious one: find $u\in\mathring H^1$ such that}
\[
\int\operatorname{grad}u\cdot\operatorname{grad}v=\int fv,\qquad v\in\mathring H^1.
\]
{The mixed weak formulation seeks $\sigma\in H(\operatorname{div})$, $u\in L^2$ such that}
\[
\begin{gathered}
\int\sigma\cdot\tau-\int u\operatorname{div}\tau=0,\qquad\tau\in H(\operatorname{div}),\\
\int\operatorname{div}\sigma\,v=\int fv,\qquad v\in L^2.
\end{gathered}
\]
{The Dirichlet condition is natural in the mixed formulation.}

% block: hodge-r3-table-summary page:58 type:ordinary-paragraph
{Table~\ref{tab:hodge-r3-boundary-conditions} collects key features of the Hodge Laplacian problems for the de Rham complex.}

% block: hodge-r3-boundary-duality-heading page:58 type:subsection
\subsubsection{{The de Rham complex with boundary conditions and duality}}\label{subsec:de-rham-boundary-duality}

% block: hodge-r3-boundary-complex page:58 type:ordinary-paragraph
{We have just considered the application of the abstract theory to the $L^2$ de Rham complex, whose domain complex is given by \eqref{eq:l2-de-rham-domain-complex}. Instead, we may apply the theory to the \emph{$L^2$ de Rham complex with boundary conditions}, by which we mean the Hilbert complex whose base spaces remain $L^2$ but whose domain complex is}
\[
0\to\mathring H^1\xrightarrow{\operatorname{grad}}\mathring H(\operatorname{curl})\xrightarrow{\operatorname{curl}}\mathring H(\operatorname{div})\xrightarrow{\operatorname{div}}L^2\to0.
\]
{This Hilbert complex is essentially the dual complex to \eqref{eq:l2-de-rham-domain-complex}, i.e., \eqref{eq:l2-dual-domain-complex}, except for some signs and for a change of indexing. Specifically, the index $k$ space of the de Rham complex with boundary conditions is the index $3-k$ space of the dual complex to the original de Rham complex. The dual complex to the de Rham complex with boundary conditions is then}
\[
0\leftarrow L^2\xleftarrow{-\operatorname{div}}H(\operatorname{div})\xleftarrow{\operatorname{curl}}H(\operatorname{curl})\xleftarrow{-\operatorname{grad}}H^1\leftarrow0,
\]
{i.e., the original de Rham complex, except for signs and indexing.}

% block: hodge-r3-boundary-table page:59 type:table
\begin{table}[htbp]\centering\small
\caption{{The Hodge Laplace problem for the de Rham complex with boundary conditions. All boundary conditions are essential for the mixed formulation and underlined boundary conditions are essential for the primal formulation. See Table~\ref{tab:hodge-r3-boundary-conditions}.}}\label{tab:hodge-r3-relative-boundary-conditions}
\setlength{\tabcolsep}{3pt}
\begin{tabular}{cccll}\hline
$k$&$L^k$&$\sigma$&{Boundary conditions}&{Harmonic forms}\\\hline
$0$&$-\Delta$&$0$&Dirichlet: $\underline{u=0}$&$\mathring{\mathsf H}^0=0$\\\noalign{\vskip4pt}
$1$&$-\vec\Delta$&$-\operatorname{div}u$&{electric}: $\underline{u\times n=0}$, $\operatorname{div}u=0$&$\mathring{\mathsf H}^1=\mathsf H^2$ $(\dim b_2)$\\\noalign{\vskip4pt}
$2$&$-\vec\Delta$&$\operatorname{curl}u$&{magnetic}: $\underline{u\cdot n=0}$, $\operatorname{curl}u\times n=0$&$\mathring{\mathsf H}^2=\mathsf H^1$ $(\dim b_1)$\\\noalign{\vskip4pt}
$3$&$-\Delta$&$-\operatorname{grad}u$&Neumann: $\partial u/\partial n=0$&$\mathring{\mathsf H}^3=\mathsf H^0$ $(\dim b_0)$\\\hline
\end{tabular}
\end{table}

% block: hodge-r3-poincare-lefschetz page:59 type:ordinary-paragraph
{If we write $\mathring{\mathcal B}^k$, $\mathring{\mathcal Z}^k$, and $\mathring{\mathsf H}^k$ for the spaces of boundaries, cycles, and harmonic forms of the complex with boundary conditions, then $\mathring{\mathcal B}^k=\mathcal B^{3-k}$ and similarly for $\mathcal Z$ and $\mathsf H$. (The relation $\mathsf H^{n-k}=\mathring{\mathsf H}^k$ on an $n$-dimensional manifold is known as \emph{Poincaré–Lefschetz duality}. For manifolds without boundary it becomes $\mathsf H^{n-k}=\mathsf H^k$, which is Poincaré duality.) It is easy to check that the strong formulation and the primal weak formulation for the complex with boundary conditions coincide with those for the ordinary de Rham complex, except for an index change, so we obtain the same boundary value problems as above. The $k=0$ problem, for instance, is the Poisson equation with Dirichlet boundary conditions. The mixed weak formulation for the problems with $k>0$, however, brings in something new. For $k=1$, for example, we seek $\sigma\in\mathring H^1$, $u\in\mathring H(\operatorname{curl})$, $p\in\mathring{\mathsf H}^1=\mathsf H^2$ such that}
\[
\begin{aligned}
\int\sigma\tau-\int u\cdot\operatorname{grad}\tau&=0,&\tau&\in\mathring H^1,\\
\int\operatorname{grad}\sigma\cdot v+\int\operatorname{curl}u\cdot\operatorname{curl}v+\int p\cdot v&=\int f\cdot v,&v&\in\mathring H(\operatorname{curl}).\\
\int u\cdot q&=0,&q&\in\mathsf H^2.
\end{aligned}
\]

% block: hodge-r3-relative-mixed-comparison page:59 type:ordinary-paragraph
{This is a mixed formulation for the vector Laplacian with boundary conditions $\operatorname{div}u=0$, $u\times n=0$. These are the same boundary conditions as for the $k=2$ problem for the ordinary de Rham complex. However, the mixed method \eqref{eq:hodge-r3-electric-mixed} for that problem was altogether different, using a vector field for the auxiliary variable $\sigma$ rather than a scalar field, as here, and imposing the boundary conditions naturally, rather than essentially as here.}

% block: hodge-r3-relative-summary page:59 type:ordinary-paragraph
{A similar situation holds for $k=2$ and $k=3$, which we leave to the reader to work out in detail. A summary is presented in Table~\ref{tab:hodge-r3-relative-boundary-conditions}.}

% block: hilbert-summary-heading page:60 type:ordinary-paragraph
\begin{tcolorbox}[enhanced,breakable,colback=KuangNei,frame hidden,boxrule=0pt,sharpish corners,before skip=10pt,after skip=10pt]
\refstepcounter{booksummary}\label{box:hilbert-complexes-summary}
\noindent\textbf{Box~\thebooksummary:} {Summary of Chapter~\ref{chap:hilbert-complexes} on Hilbert complexes.}\par

% block: hilbert-summary-definition page:60 type:ordinary-paragraph
{A \emph{Hilbert complex} is a sequence of Hilbert spaces $W^k$ and closed densely defined operators $d^k$ from one to the next such that $d^k$ maps its domain into the null space of $d^{k+1}$. The Hilbert complex is called \emph{closed} if the $d^k$ all have closed range.}

% block: hilbert-summary-domain-complex page:60 type:ordinary-paragraph
{Setting $V^k=D(d^k)$, we get the domain complex}
\[
0\to V^0\xrightarrow{d}V^1\xrightarrow{d}\cdots\xrightarrow{d}V^n\to0,
\]
{a cochain complex of Hilbert spaces and bounded linear operators.}

% block: hilbert-summary-cocycles page:60 type:ordinary-paragraph
{We denote the cocycles and coboundaries as $\mathcal Z^k$ and $\mathcal B^k$ and, for the dual complex, by $\mathcal Z_k^*$ and $\mathcal B_k^*$.}

% block: hilbert-summary-laplacian page:60 type:ordinary-paragraph
{The \emph{Hodge Laplacian} is the unbounded operator $L^k$ from $W^k\to W^k$ given by $d^*d+dd^*$ with the natural domain. Its kernel is the space $\mathsf H^k=\mathcal Z^k\cap\mathcal Z_k^*$ of \emph{harmonic forms}. For a closed Hilbert complex, $\mathsf H^k$ is isomorphic to the cohomology space $\mathcal H^k$.}

% block: hilbert-summary-decomposition page:60 type:ordinary-paragraph
{Any closed Hilbert complex gives rise to the orthogonal \emph{Hodge decomposition}}
\[
W^k=\mathcal B^k\BookOrthSum\mathsf H^k\BookOrthSum\mathcal Z^k
\]
{and satisfies the \emph{Poincaré inequality}}
\[
\|z\|_V\leq c_P\|dz\|,\qquad z\in\mathcal Z^{k\perp_V}.
\]

% block: hilbert-summary-solvability page:60 type:ordinary-paragraph
{The Hodge Laplace problem}
\[
L^ku=f-P_{\mathsf H}f,\qquad u\perp\mathsf H^k,
\]
{has a unique solution in $D(L^k)$ for any $f\in W^k$. The solution may be characterized as well by the primal weak formulation \eqref{eq:hodge-primal-formulation} or the mixed weak formulation \eqref{eq:hodge-mixed-formulation}.}
\end{tcolorbox}
